\documentclass[10pt,aspectratio=43,mathserif]{beamer}		
%设置为 Beamer 文档类型，设置字体为 10pt，长宽比为16:9，数学字体为 serif 风格

%%%%-----导入宏包-----%%%%
\usepackage{ecnu}			%导入 CCNU 模板宏包
\usepackage{ctex}			 %导入 ctex 宏包，添加中文支持
\usepackage{amsmath,amsfonts,amssymb,bm}   %导入数学公式所需宏包
\usepackage{color}			 %字体颜色支持
\usepackage{graphicx,hyperref,url}	
\usepackage{minted,fontspec}
%%%%%%%%%%%%%%%%%%	

\setmonofont{Source Code Pro}
\definecolor{bg}{rgb}{0.95,0.95,0.95}
\newminted{asm}{bgcolor=bg}
\newminted{cpp}{bgcolor=bg}


\beamertemplateballitem		%设置 Beamer 主题

%%%%------------------------%%%%%
\catcode`\。=\active         %或者=13
\newcommand{。}{．}				
%将正文中的“。”号转换为“.”。
%%%%%%%%%%%%%%%%%%%%%

%%%%----首页信息设置----%%%%
\title[ECNUBeamerTemplate]{并行计算大作业}
\subtitle{——并行优化矩阵乘法}			
%%%%----标题设置


\author[x]{
%   ~ \\\medskip
  {\small {朱桐~10175102111}} \\
  {\small {周亦然~10175102207}}
  }
%%%%----个人信息设置

% \institute[ECNU]{
%   East China Normal University \\
%   Department of Computer Science and Technology}
%%%%----机构信息

\date[2019.12.27]{
  2019.12.27}
%%%%----日期信息

\begin{document}

\begin{frame}
\titlepage
\end{frame}				%生成标题页

\section{Outline}
\begin{frame}
\frametitle{Outline}
\tableofcontents
\end{frame}				%生成提纲页


\section{Introduction}
\begin{frame}
  \frametitle{Backgrounds}
  \begin{itemize}
    \item {1111111111111}
    \item {2222222222222}
    \item {3333333333333}
    \item {4444444444444}
  \end{itemize}
\end{frame}

\section{Photo}
\begin{frame}
  \frametitle{My Photo}
    \begin{figure}[htbp]
        \centering
        \includegraphics[width=0.6\textwidth]{pic.png}
        \caption{hahahaha...}\label{fig:digit}
    \end{figure}
\end{frame}

\section{Conclusion}
\begin{frame}
  \frametitle{Emmm...}
  \noindent\textbf{Sequence Tagging Loss\\}
    \[{\mathcal{L}_p} =  - \sum\limits_{i = 1}^S {\sum\limits_{j = 1}^N {{p_{i,j}}\log ({{\hat p}_{i,j}})} } \]
  \noindent\textbf{Language Classifier Loss\\}
   \[{\mathcal{L}_a} =  - \sum\limits_{i = 1}^S {{l_i}\log ({{\hat l}_i})}\]
  \noindent\textbf{Bidirectional Language Model Loss\\}
   \[{\mathcal{L}_l} =  - \sum\limits_{i = 1}^S {\sum\limits_{j = 1}^N {\log (P({w_{j + 1}}|{f_j})) + \log (P({w_{j - 1}}|{b_j}))} }\]
\end{frame}
 
\section{References}
\begin{frame}{References}
\begin{thebibliography}{99}
\bibitem{zhao1} Xuezhe Ma and Eduard Hovy. (2016).\\
{\bf End-to-end Sequence Labeling via Bi-directional LSTM-CNNs-CRF.\\}
In Proceedings of the 54th Annual Meeting of the Association for Computational Linguistics, pages 1064–1074, Berlin, Germany, August 7-12, 2016.
\bibitem{zhao2} Marek Rei. (2017).\\
{\bf Semi-supervised Multitask Learning for Sequence Labeling.\\}
In Proceedings of the 55th Annual Meeting of the Association for Computational Linguistics, pages 2121–2130, Vancouver, Canada, July 30 - August 4, 2017.
\end{thebibliography}
\end{frame}

\section{CUDA}


\begin{frame}
    \frametitle{solving matrix mult on CUDA}

    \begin{itemize}
        \item host vs device, distributed memory
        \item streamprocessors
        \item blocks, wraps, threads, shared memory
    \end{itemize}

    子任务：给定 $i,j$, 计算 $C_{ij}$，也就是并行最外面的两个for
\end{frame}

\begin{frame}
    \frametitle{solving matrix mult on CUDA}

    \begin{itemize}
        \item 轻量级线程，切换成本较低
        \item 基本可以平等划分任务
        \item 一个线程对应一个子任务（开 $n \times m$ 个线程）
    \end{itemize}
\end{frame}

\begin{frame}
    \frametitle{shared memory}
    \begin{itemize}
        \item 公式: $C_{ij} = \sum\limits_{k=0}^{bm} A_{ik} \cdot B_{kj}$
        \item 对于内层的两个循环$j,k$，发现公用 $A_{ix}$
        \item 使用 share memory 缓存 $A$
    \end{itemize}
\end{frame}

\begin{frame}
    \frametitle{share memory}
    \begin{itemize}
        \item block内公用
        \item \texttt{block\_size} 必须是 $am$ 的因子，保证同块内使用相同的一行
        \item 别忘了 \texttt{\_\_sychronize}
    \end{itemize}
\end{frame}

\begin{frame}
    \frametitle{Fast IO: buffered input \& output}
    \begin{itemize}
        \item \texttt{fread, fwrite instead of printf, scanf}
        \item no need to flush every time we read/write a number
    \end{itemize}
\end{frame}

\begin{frame}
    \frametitle{Occupancy: optimal size for the number of blocks and threads}
    \begin{itemize}
        \item 以 wrap 为计算单元
        \item thread 中有register
        \item block 中有share memory
        \item 每个 multiprocessor 也有最大的 register 和 share memory 和  active threads
        \item 如果资源不够将无法派出所有的资源进行运算
    \end{itemize}
\end{frame}

\begin{frame}
    \frametitle{Occupancy: optimal size for the number of blocks and threads}
    \begin{itemize}
        \item 我们需要让最大比例的显卡工作
        \item 引入 Occupancy 机制
        \item 通过 nvidia 官网上的 Occupancy calculator 进行运算
        \item \texttt{nvcc --ptxas-options=-v to nvcc} 查看内核函数使用寄存器个数
        \item 使用 \texttt{}
    \end{itemize}
\end{frame}


\begin{frame}
    \frametitle{asychronize IO}
\end{frame}

\end{document}
